\section{Native coefficients: a completed controller and a smaller radix}\label{sec:native}

The direct quartic is the principal result of this report. A separate
native route keeps fifteen positive integer variables and a 66-operation
polynomial of degree 157, while reducing the size of its four fixed
integer coefficients. This section separates a completed prior proposal
from a new change to the encoding.

\subsection{The polynomial is fixed before the controller is chosen}

The fifteen independent positive variables are
\[
 J,F,\alpha,z,f,h,i,T,s,w,\tau,\eta,\zeta,y,\gamma.
\]
For a specified whole coefficient tuple $(\beta,K_0,M_C,M_F)$, define
the following polynomial abbreviations:
\[
\begin{aligned}
 q&=\beta J+1,& C&=q-F-\alpha,\\
 X&=wq,& Y&=sq^3,\\
 E&=XY,& k&=\eta+\zeta,\\
 c&=kY+\eta,& \ell&=EkY,\\
 a&=E+Y,& H&=4a+3,\\
 \Delta&=a^2+H,& D&=X+ac+\gamma H,\\
 V&=f^2(cT-R)-c,& S&=i\Delta c^2.
\end{aligned}
\]
\[
 R=[(q-1)(q-F)+\alpha](q^2-1)+(M_C+qM_F)J.
\]
Put
\begin{align*}
 N_0&=\tau^2-\ell(\ell+k),&N_1&=D^2-\Delta c^2,\\
 N_a&=S^2(V^2-y^2)+y^2,&N_k&=k-hE-R,\\
 N_t&=(K_0+w)C+(q-F)-z(q-1),&N_s&=\Delta f^2-S^2.
\end{align*}
The single equation is
\begin{equation}\label{eq:p66}
 \boxed{N_0N_1N_aN_kN_tN_s-\Delta=0.}
\end{equation}
Its exact straight-line circuit has 38 multiplications, 15 additions,
and 13 subtractions. The apparent top terms of $N_1$ cancel: with $Z=X+\gamma H$,
$N_1=2acZ+Z^2-Hc^2$. The six factor degrees are therefore
$22,18,62,7,2,46$. Setting all inputs equal to an indeterminate $t$
gives leading term $2^{17}\beta^{91}t^{157}$, so the degree is exactly
157. The four named constants are not quantified unknowns. The controller prime below is not part of the coefficient
tuple; it is an explicit witness used in the completeness proof.

\subsection{Closing the conditional radix-$2^{49}$ construction}

The repository's shared-marker radix49 proposal~\cite{radix49} already
contains the encoding and its ordinary native/Pell proof. Its remaining
conditions were finite controller conditions. We complete precisely
those conditions and materialize the entire coherent tuple.

Let $\rho=2^{49}$, $L=8807$, $B=\rho^L$, and $d=49L=431543$.
Choose the fixed prime
\[
 \lambda=434011>d,\qquad \lambda-1=1887\cdot230.
\]
For $H_{\rm ctrl}=1+BK_0+B^2$, exact computation gives:
\begin{center}
\begin{tabular}{@{}rrrr@{}}
\toprule
Modulus & $B$ & $K_0$ & $H_{\rm ctrl}$\\
\midrule
7 & 4 & 5 & 2\\
49 & 18 & 47 & 44\\
8807 & 1901 & 973 & 3135\\
434011 & 398487 & 335571 & 6522\\
\bottomrule
\end{tabular}
\end{center}
The exact orders are
\[
 \operatorname{ord}_{49}(B)=3,\quad
 \operatorname{ord}_{8807}(B)=629,\quad
 \operatorname{ord}_{434011}(B)=18870.
\]
The required valuation-one lift is
\[
 B^{\lambda-1}\equiv96356084144
 =1+222013\lambda\pmod{\lambda^2}.
\]
Since $0<222013<\lambda$, the lift has valuation exactly one.
The full modulus 49 is checked explicitly, rather than only its prime
divisor seven. Deterministic trial division proves the asserted primes;
annihilating powers and every maximal proper-power test prove the orders.

These facts discharge all the finite premises of the inherited
conditional theorem. Its carry, phase-alignment, safe-dummy, rank-tuning,
Pell, and refutation arguments then give the ZFC-provable equivalence
between positive unsolvability of~\eqref{eq:p66} at this tuple and
$\Con(\ZFC)$. The new work here is the finite completion and complete
coefficient specialization; the shared-marker architecture is prior work.

The materialized radix49 tuple has bit lengths
$(431543,433602,431543,431544)$ in the port order
$(\beta,K_0,M_C,M_F)$, totaling 216,030 raw unsigned bytes. Its four
deterministic gzip files total 8,774 bytes. The corresponding matched
radix53 and radix55 tuples occupy 233,666 and 242,484 raw bytes.
Raw coefficient storage, compressed storage, and printed-page length
are different metrics. The complete coefficient recipe, binary files,
standalone verifier, and inherited proof texts are supplied under
\file{radix49/}.

\subsection{Complemented labels permit a radix-$2^{47}$ encoding}

The radix49 completion suggests a further change to the coefficient
construction. The previous radix47 local obstruction~\cite{radix47obstruction}
concerned unchanged selector coefficients: its predecoded carry allowance
was too large to infer that a cell selected at most one tile. Instead of
discarding that obligation, change the numerical labels so the full
coefficient sum becomes smaller.

For each original tile $(w_i,n_i,p_i,s_i)$, apply the bijections
\[
 w_i'=255-w_i,\quad p_i'=255-p_i,\quad
 n_i'=511-n_i,\quad s_i'=511-s_i.
\]
The west/east labels use the same eight-bit bijection and the north/south
labels the same nine-bit bijection. Thus edge equality and inequality,
including the typed Start and forbidden-color conditions, are preserved.
The seven physical marker bits and the occupancy bits are kept. An empty
cell still has zero physical label bits; complementation is the relabeling
of genuine tiles, not an unconditional complement of every empty bit.

Put
\[
 u_i=w_i+2^8n_i+2^{17}p_i+2^{25}s_i.
\]
The complemented packed label is $2^{34}-1-u_i$, so the new selector
coefficient is
\begin{equation}\label{eq:complement}
 c_i=1+2^{10}(u_i+1)\qquad(0\le i\le942).
\end{equation}
Every nonselector coefficient and physical support position is as in the
shared-marker construction. Exact reconstruction of all 943 tiles gives
\[
 \sum_e c_e=3,953,059,820,717,999,
 \qquad\max_i c_i=15,049,678,914,561<2^{44}.
\]
The old complete coefficient sum was 12,917,846,596,842,415. Both sums
include the nonselector controls; the reduction is not obtained by
omitting any coefficient from the proof.

\subsubsection*{The full new coefficient tuple}

Now set $\rho=2^{47}$, $L=8807$, $B=\rho^L=2^{413929}$,
$\beta=B-1$, and $\Lambda=\rho^{4393}+\rho^{5495}$.
Besides~\eqref{eq:complement}, the complete nonzero controls are
\begin{align*}
 c_{1023+j}&=2^{10+j},&c_{1003+j}&=2^{27+j}&&(0\le j\le7),\\
 c_{1013+j}&=2^{18+j},&c_{983+j}&=2^{35+j}&&(0\le j\le8),\\
 c_e&=2^{44}&&&
 e&\in\{992,1011,1022,1031,1095,2195,3296\}.
\end{align*}
All other $c_e$ vanish. Let $\mathcal E$ be their support together with
the two coefficient-zero dummy positions 1063 and 1064. It has 986
positions. Define
\begin{align*}
 K_0={}&\rho^{30}+\rho^{62}+\rho^{1100}+\rho^{2201}+\rho^{8849}
       +\sum_{e=0}^{3296}c_e(\rho^{4393-e}+\rho^{5495-e}),\\
 M_C={}&\beta-\sum_{e\in\mathcal E\setminus\{1\}}\rho^e-2\rho^{1063},\\
 C_{\rm core}={}&\sum_{e=1013}^{1022}\rho^e+
              \sum_{e=1065}^{1073}\rho^e+\rho^{2195}+\rho^{3296},\\
 M_F={}&2B+3+(2^{47}-2)\Lambda+C_{\rm core}-\rho^{8792}
                +(2^{28}-3)\rho^{7694}-\rho^{7691}.
\end{align*}
These are fixed integers. Every exponent and finite sum limit is fixed.
The exact 943-tile data and the complete 3297-slot coefficient array are
supplied; all four binary integers are also materialized. In particular,
no unspecified tile enumeration or proof search is needed to define them.

The positive native high-mask padding behind the compressed $M_F$ formula
is fifteen full fields at 8792--8806, three full fields at 7691--7693,
and bits 2--27 at field 7694. Its population is
$15\cdot47+3\cdot47+26=872$. Together with the unchanged 114 core
tests, this gives 986 ones, exactly matching the physical low support.
The source shift by $\beta$ is applied once. All four ports are changed
coherently; none is imported numerically from the radix49 tuple.

\subsubsection*{Why the old carry obstruction no longer applies}

The complete predecoded convolution bound, including the exceptional
dummy and every temporal phase, maps the carry interval $0,\ldots,29$
into itself. Thus
\[
 \#\{\hbox{selected tiles}\}+k\le943+29=972<1024.
\]
The ten low tested bits therefore recover zero or one selected tile.
The remaining payload bits recover the complemented label and, by the
same shared-marker counting argument, all seven marker bits individually.

After this decoding, the selector maximum is below $\rho/8$ and the
entire nonselector control sum is $\rho-1024$. The complete decoded
convolution and temporal estimate, even with the earlier incoming carry,
is strictly below $13\rho$. Hence the actual outgoing carries are at
most 12. This bound protects the anchor predecessors and all subsequent
local resets.

At radix47 an occupied clause has raw value $\rho+1$. Its canonical
digit is one and it emits carry one. This carry must not be silently
dropped. The proof explicitly resets it before every later tested copy,
anchor, padding, or origin field. For example, the larger clause
predecessor has raw bound at most $1+7\cdot2^{44}-2^{10}$, and adding
12 still leaves it below $\rho$. The two phase-anchor predecessor
identities also leave room for every temporal phase plus carry 12.
They recover integral temporal phase and whole-cell alignment as before.
The raw padding predecessors at 7690 and 8791 are bounded by a single
selector coefficient plus the incoming carry, so both reset. The complete
zero fields 7691--7695 and 8792--8806 then validate the new padding.

These arguments preserve the positive-source slack, all simultaneous
safe-dummy changes, and the actual rank identity. They supply the same
canonical $R\equiv3\pmod4$ and
$\pc(R)=3dN+2$ interface required by the P66 converse. The full ordinary
proof and the independent review are included in \file{radix47/}.

\subsubsection*{The new finite controller}

Here $d=47\cdot8807=413929$. Choose
\[
 \lambda=1012691>d,\qquad
 \lambda-1=10\cdot101269,
 \qquad101269=\operatorname{lcm}(23,4403).
\]
The exact residues are:
\begin{center}
\begin{tabular}{@{}rrrr@{}}
\toprule
Modulus & $B$ & $K_0$ & $1+BK_0+B^2$\\
\midrule
47 & 12 & 5 & 17\\
8807 & 2677 & 4079 & 5042\\
1012691 & 959313 & 1005855 & 834350\\
\bottomrule
\end{tabular}
\end{center}
The exact orders at these three moduli are respectively
$23$, $4403$, and $1012690$. Finally,
\[
 B^{\lambda-1}\equiv1+282338\lambda\pmod{\lambda^2}.
\]
All controller units are nonzero and the lift has valuation one. The
same cofinal CRT and Boolean-subset rank-tuning proof now applies with
the new $d$, prime, spacing, and tuple. Together with the complete
carry/alignment proof, it establishes the ZFC-provable equivalence
between positive unsolvability of~\eqref{eq:p66} at this tuple and
$\Con(\ZFC)$.

\begin{center}
\begin{tabular}{@{}lrrrrr@{}}
\toprule
Tuple & $\beta$ bits & $K_0$ bits & $M_C$ bits & $M_F$ bits & Raw bytes\\
\midrule
Radix55 & 484385 & 486696 & 484385 & 484386 & 242484\\
Radix53 & 466771 & 468998 & 466771 & 466772 & 233666\\
Completed radix49 & 431543 & 433602 & 431543 & 431544 & 216030\\
New radix47 & 413929 & 415904 & 413929 & 413930 & \textbf{207214}\\
\bottomrule
\end{tabular}
\end{center}

The new tuple saves 26,452 raw bytes, about 11.32\%, against the verified
radix53 tuple. Its compressed four-coefficient payload is 8,356 bytes.
The scalar $B=2^{413929}$ is also smaller than the collection's previously
stated $2^{429959}$ native-base endpoint; this statement concerns that
scalar, not a comparison of every metric of the two different encodings.

Label complementation itself has predecessors in the repository. The new
result is this complete indexed943/shared-marker/radix47 composition,
the corrected carry bounds and resets, its entire coherent tuple, and the
completed controller. It is not a claim to have invented complementation
or to have proved a globally smallest coefficient tuple.

